\protect \section *{Розділ~\ref *{sec:mhd}. \nameref *{sec:mhd}}
\begin{Solution}{5.{1}}
		$p=B^2/8\pi=2\pi n^2I^2/c^2$, $B=4\pi nI/c$.
	
\end{Solution}
\begin{Solution}{5.{2}}
		$p=2nT=3{,}2\times10^{6}$~дин/см$^2$, $B(\beta=1)\approx9\times10^{3}$~Гс; $\beta\approx3{,}2\,\%$ за $5\times10^4$~Гс.
	
\end{Solution}
\begin{Solution}{5.{3}}
		$\sigma\approx3{,}6\times10^{17}~\mathrm{с}^{-1}$ (порівнянна з~міддю), $\tau_d=4\pi\sigma L^2/c^2\approx50$~с, $R_m\approx5\times10^4$.
	
\end{Solution}
\begin{Solution}{5.{4}}
		$v_A\approx2{,}2\times10^{8}$~см/с, $v_{Ti}\approx1{,}0\times10^{8}$~см/с ($v_{Ti}/v_A\approx0{,}45$).
	
\end{Solution}
\begin{Solution}{5.{5}}
		$B^2/8\pi\approx1{,}1\times10^{8}$~дин/см$^2\approx112$~бар, $\beta\approx3\,\%$.
	
\end{Solution}
\begin{Solution}{5.{6}}
		$v_A\approx7{,}7\times10^{8}$~см/с, $t=L/v_A\approx1{,}3\times10^{-7}$~с.
	
\end{Solution}
\begin{Solution}{5.{7}}
		$\lambda_{\mathrm{пр}}\sim10^{6}$~см (10~км), $r_{Li}\approx0{,}3$~см, для іонів $D_\perp/D_\parallel\sim(\nu_{ii}/\omega_{ci})^2\sim10^{-13}$, для електронів $(\nu_{ei}/\omega_{ce})^2\sim10^{-17}$.
	Уздовж поля частинки вільно пролітають $\lambda_{\mathrm{пр}}\gg L$, а~поперек поля гірація з~$r_L\ll L$ локалізує їх: роль довжини пробігу відіграє $r_L$.
	
\end{Solution}
\begin{Solution}{5.{8}}
		$\Gamma_y=c\,T\,n^\prime/(qB)$, $j_y=c\,p^\prime/B$.
	
\end{Solution}
\begin{Solution}{5.{9}}
		у~$10^4$ разів.
	
\end{Solution}
\begin{Solution}{5.{10}}
		$\psi\approx45^\circ$.
	
\end{Solution}
\begin{Solution}{5.{11}}
		$W=\tfrac34RB_0|\Delta\Phi|\approx1{,}7\times10^{13}$~ерг~$=1{,}7$~МДж; ні, петля дає лише $\int p\,dS$, тобто добуток~$nT$.
	
\end{Solution}
\begin{Solution}{5.{12}}
		$\dfrac{4\pi n_cS}{c^2}\,I\approx4\times10^{-4}$~СГС$\,=0{,}13$~В$\cdot$с (для $I=10^6$~А~$=3\times10^{15}$~СГС);
не залежить~--- лише від повного охопленого струму.
	
\end{Solution}
\begin{Solution}{5.{13}}
		(а)~$\eta_\parallel=5{,}2\times10^{-5}\cdot15/1000^{3/2}\approx2{,}5\times10^{-8}$~Ом$\cdot$м
		(близько до~міді, $1{,}7\times10^{-8}$~Ом$\cdot$м), $\eta_\perp=2\eta_\parallel\approx4{,}9\times10^{-8}$~Ом$\cdot$м.
		(б)~$\nu_{ei}=2{,}91\times10^{-6}\cdot10^{14}\cdot15/1000^{3/2}\approx1{,}4\times10^{5}$~с$^{-1}$,
		$\omega_{ce}=1{,}76\times10^{7}\cdot5\times10^{4}\approx8{,}8\times10^{11}$~с$^{-1}$, тож
		$(\omega_{ce}/\nu_{ei})^2\approx4\times10^{13}$ і~$\eta_\perp\sim10^{6}$~Ом$\cdot$м~--- на~13 порядків більше за~реальний; плазма
		була б практично непровідною поперек поля, що суперечить спостереженню.
		(в)~Тертя $\vect R_{ei}=en\vect j/\sigma$ залежить лише від відносної швидкості електронів та~іонів, тобто від~$\vect j$.
		Спільний дрейф $\vect E\times\vect B$ однаковий для обох сортів, відносної швидкості не створює й~тертя не викликає;
		опір поперек поля визначає лише доданок $\vect j_\perp$, тому $\eta_\perp\sim\eta_\parallel$.
	
\end{Solution}
